\part{The Exercises}

\chapter{Constrained Literature}
\label{ch:constrained}

A line such as \texttt{a sad lad asks dad;} is language generated under an expanding admissibility boundary. This chapter treats such lines as a genre.

\section{The Phenomenon}

If $\Kset_i$ is the current key set and $\Sig(w)$ the characters a word $w$ requires, the vocabulary available to a lesson author is
\[
\Vocab_i=\{\,w\in\Vocab:\Sig(w)\subseteq\Kset_i\,\}.
\]
The author does not choose freely from English but from $\Vocab_i$, and the peculiar texture of early exercises is the visible shape of that restriction. In the first exercises of the Academy pages the restriction produces \texttt{ass lad add alf} and \texttt{flak alas fall asks} in Exercise~1 and the phrase \texttt{a sad lad asks dad;} that closes Exercise~2. In Exercise~2 it produces \texttt{gaff gash lags hall} and the marginal words \texttt{gald} and \texttt{jaffa}.

\section{Two Regimes}

Chapter~8 distinguished two regimes. In the \emph{admissibility regime} the constraint is what can be spelled, and the author strains for words. In the \emph{salience regime} the constraint is what will exercise the new key,
\[
\Vocab_i^{\mathrm{new}}=\{\,w\in\Vocab_i:\Sig(w)\cap\Newk_i\neq\varnothing\,\}.
\]
Exercise~13 is entirely in the salience regime: \texttt{zip zoo zig zag}, \texttt{zinc: zulu zero jazz:}, \texttt{quiz doze size hazy:}, \texttt{zeal lazy fizz fuzz:}. The sentences on the pages that continue Exercises 6 and later show the same regime at the level of the sentence: \texttt{Ilse who is a redhead is a selfish girl who likes a wild life;} is a sentence chosen so that every content word contains W or I, and its meaning is incidental.

\section{Kinship with Constrained Writing}

The genre is a relative of the lipogram, a text written without one or more letters. Georges Perec's novel \emph{La Disparition} (1969) avoids the letter E throughout, and the lessons here are the pedagogical inverse: they permit a small set and exclude the rest. The comparison is instructive because in each case the constraint is announced by the format and its effects are read in the prose. In a typing lesson the constraint is announced by the diagram at the top of the page.

\section{The Successive Alphabet}

The following passages are reconstructed, not transcribed. They show how a text can be written under a sequence of enlarged key sets. The first line is observed; the rest are constructed.

\begin{Verbatim}[label=QWERTY, labelposition=topline]
a sad lad asks dad;                        (observed, keys: A S D F J K L ;)
a glad lad has a flask;                    (constructed, adds G H)
a sad lad has used a glass flask;          (constructed, adds E U)
a lad has a rose; dad asks for a jug;      (constructed, adds R O)
today a lad dared to ask dad;              (constructed, adds T Y)
\end{Verbatim}

The prose begins almost as incantation, grows faintly surreal, and resolves into ordinary language only when the alphabet nearly completes.

\section{Two Constraints Running in Opposite Directions}

An acrostic mnemonic restricts the initial letters of the words in a phrase. A restricted-alphabet lesson restricts all the letters. The first constraint runs from a target set to a phrase; the second runs from a permitted set to a vocabulary. Chapter~24 shows self-devised mnemonics of the first kind for the Dvorak keyboard, and the contrast clarifies what each constraint is for: the acrostic fixes a position in memory, and the lesson fixes a movement in the hand.

\section{Measuring a Restriction}

The size $|\Vocab_i|$ is a measurable property of a lesson. Appendix~D gives the procedure, and Table~\ref{tab:vocab} reports the result for the design key order of Chapter~\ref{ch:qwerty}, for QWERTY and for Dvorak by position matching. Letter matching has the QWERTY column by definition. Counts are taken from two lists. The first is a cleaned copy of the google-10000-english list (9\,486 entries after removing short tokens, place names, and brands), and the column ``cov.'' is the share of rank-weighted frequency (weight $1/\mathrm{rank}$) that the admissible words capture. The second, used only in the text below, is a 370\,105-entry list of English words that includes many rare and inflected forms.

\begin{table}[ht]
\centering
\caption{Admissible vocabulary by exercise under the design key order. Letters only; punctuation keys are omitted from the ``new'' columns because they form no words.}
\label{tab:vocab}
\begin{center}\small
\begin{tabular}{@{}rlrrrlrrr@{}}
\toprule
& \multicolumn{4}{c}{QWERTY (and letter matching)} & \multicolumn{4}{c}{Dvorak by position} \\
\cmidrule(lr){2-5}\cmidrule(lr){6-9}
Ex. & new & $|\Vocab_i|$ & $|\Vocab_i^{\mathrm{new}}|$ & cov.\ \% & new & $|\Vocab_i|$ & $|\Vocab_i^{\mathrm{new}}|$ & cov.\ \% \\
\midrule
1 & \texttt{asdfjkl} & 21 & 21 & 1.1 & \texttt{aehnostu} & 185 & 185 & 20.7 \\
2 & \texttt{gh} & 36 & 15 & 1.5 & \texttt{di} & 477 & 292 & 32.1 \\
3 & \texttt{eu} & 175 & 139 & 3.3 & \texttt{g} & 670 & 193 & 33.4 \\
4 & \texttt{ro} & 524 & 349 & 14.7 & \texttt{pr} & 1840 & 1170 & 42.2 \\
5 & \texttt{ty} & 1102 & 578 & 35.8 & \texttt{fy} & 2302 & 462 & 54.1 \\
6 & \texttt{wi} & 1819 & 717 & 46.9 & \texttt{c} & 3418 & 1116 & 60.4 \\
7 & \texttt{nm} & 4089 & 2270 & 70.6 & \texttt{bm} & 4951 & 1533 & 72.7 \\
8 & \texttt{vb} & 5360 & 1271 & 79.5 & \texttt{kx} & 5471 & 520 & 75.1 \\
9 & \texttt{c} & 7366 & 2006 & 89.8 & \texttt{jw} & 6009 & 538 & 81.6 \\
10 & \texttt{qp} & 9139 & 1773 & 98.9 & \texttt{l} & 8463 & 2454 & 94.7 \\
11 & \texttt{x} & 9370 & 231 & 99.7 & \texttt{qv} & 9370 & 907 & 99.7 \\
12 & \texttt{--} & 9370 & 0 & 99.7 & \texttt{z} & 9486 & 116 & 100.0 \\
13 & \texttt{z} & 9486 & 116 & 100.0 & \texttt{--} & 9486 & 0 & 100.0 \\
\bottomrule
\end{tabular}
\end{center}
\end{table}

Four findings follow, each subject to the caveat that the frequency list is drawn from web text and the key order is a design.

\begin{enumerate}
\item \textbf{The first exercise.} Only 21 entries of the cleaned frequency list are spellable from the QWERTY home-row letters, and some of them are fragments (for example \texttt{dsl}); the same list yields 185 words for the Dvorak home-row set, including \texttt{the}, \texttt{to}, \texttt{on}, \texttt{that} and \texttt{not}. The 370\,105-entry list gives 173 against 1\,599, a ratio of about nine in both lists.
\item \textbf{Coverage.} The admissible words capture about 1.1\% of rank-weighted frequency at QWERTY Exercise~1 and about 20.7\% at the Dvorak counterpart. QWERTY does not reach a third of the weighted frequency until Exercise~5, while Dvorak by position exceeds it by Exercise~3.
\item \textbf{The advantage is early.} From Exercise~7 onward the two vocabularies are close, and at Exercises 8 to 10 QWERTY under the design order is ahead in coverage, because the design introduces N, M, V, B and C while position matching sends Dvorak to its bottom-row letters K, X, J and W. The effect belongs to the first six exercises.
\item \textbf{Punctuation positions.} Exercise~12 introduces no letter on QWERTY and Exercise~13 introduces none on Dvorak. In each case $\Vocab_i^{\mathrm{new}}$ is empty, so the salience regime cannot operate through words.
\end{enumerate}

The expectation that the first-exercise vocabulary would be far larger for Dvorak is therefore supported for the first six exercises and reversed later in this particular order.

\chapter{The QWERTY Curriculum}
\label{ch:qwerty}

This chapter states a QWERTY curriculum as a template instantiation. It is a design, not a report of the historical order, except where the two observed exercises fix an anchor. \design{the order below is proposed; only Exercises 1, 2, 6 and 13 are anchored to the Academy pages}

\section{The Ladder Template}

For an exercise $i$ with new keys $\Newk_i$ and admitted set $\Kset_i$, the template produces groups as follows.

\begin{enumerate}[label=\textbf{G\arabic*.}]
\item A chain of the keys in $\Newk_i$ (Exercises 1 and 2), or, for later exercises, a sandwich drill in which each new key is placed between two presses of the home key of its finger column (\texttt{ded}, \texttt{juj}).
\item The chain with one old key inserted (Exercises 1 and 2), or the sandwich drill with the new letter capitalized, so that shift coordination is practiced from the first reach.
\item A scramble of the first group, or short words from $\Vocab_i^{\mathrm{new}}$.
\item Words from $\Vocab_i^{\mathrm{new}}$.
\item Longer words from $\Vocab_i^{\mathrm{new}}$.
\item A phrase.
\item[\textbf{G7+.}] Sentences, in which each content word contains a new key where possible.
\end{enumerate}

Each group is printed as three lines, each a window of fixed width on the endless repetition of a unit. The template has three parameters: the key-introduction order, the lexicon, and the layout header. Appendix~G contains the corpora generated from the template for both layouts.

\section{The Key Order}

The order below is designed so that the four anchor exercises are respected and every letter is introduced by Exercise~13.

\begin{center}
\begin{tabular}{@{}cll@{}}
\toprule
Ex. & New keys & Status \\
\midrule
1 & A S D F ; L K J & anchored \\
2 & G H & anchored \\
3 & E U & design \\
4 & R O & design \\
5 & T Y & design \\
6 & W I & anchored \\
7 & N M & design \\
8 & V B & design \\
9 & C , & design \\
10 & Q P & design \\
11 & X . & design \\
12 & / and shifted punctuation & design \\
13 & Z & anchored \\
\bottomrule
\end{tabular}
\end{center}

\section{Home Row, Reaches, Alternation}

The first two exercises complete the home row. Reaches follow as sandwich drills of the form home, reach, home, exactly as in both paper sources (\texttt{sws}, \texttt{kik}, \texttt{aqa}). Alternation exercises then use words in which consecutive letters fall on opposite hands, and same-hand exercises use words typed entirely by one hand. The workbook's later lessons treat alternate-hand and one-hand words as topics in their own right, which supports treating hand alternation as a curricular category and not merely as an incidental property of words.

\section{Digraphs, Punctuation, Capitals, Numbers}

After the letters, the curriculum introduces frequent digraphs as units, then punctuation as keystrokes within sentences, then capitals through case-alternating drills, then the number row, and finally longer passages. The workbook places numbers and symbols in separate drill blocks (Lessons 6 to 10 and 11 to 15) and does so after the alphabet is complete, a sequencing that the design here follows.

\section{Longer Passages}

The final stage moves to continuous prose and to timed passages. For the timed passages the design adopts the scoring conventions found in the workbook: five-minute timings repeated twice, cumulative word counts printed at the right margin, and a policy of not correcting errors. Appendix~B states the scoring rules.

\chapter{The Dvorak Curriculum}
\label{ch:dvorakcurr}

The same ladder can be applied to Dvorak, and there are two ways to do so. This chapter derives both from the QWERTY exercises and shows what each costs.

\section{Position Matching}

Position matching transposes the QWERTY key order by physical location, using the mapping of Chapter~12. Exercise~1 becomes the keys A O E U and S N T H. The chain groups transpose cleanly:

\begin{Verbatim}
G1   aoeu snth aoeu snth ...
G2   aoeua snths aoeua snths ...
G3   eauo tshn eauo tshn ...
\end{Verbatim}

(derived by mapping the QWERTY lines \texttt{asdf ;lkj}, \texttt{asdfa ;lkj;}, \texttt{dafs k;jl}; not observed.) The word groups do not transpose. The QWERTY words \texttt{ass lad add alf} map to \texttt{aoo nae aee anu}, which are not words. A position-matched lesson can keep the geometry and must therefore replace the lexicon.

\section{Dvorak-Native Vocabulary}

Regenerating the words from $\Vocab_1$ for the Dvorak key set $\{a,o,e,u,h,t,n,s\}$ gives real and frequent words, among them \texttt{the}, \texttt{that}, \texttt{than}, \texttt{then}, \texttt{thus}, \texttt{east}, \texttt{aunt}, \texttt{haunt}, \texttt{those}, \texttt{these}, \texttt{south}, \texttt{house}, \texttt{hasten}, \texttt{tenants} and \texttt{sunset}. Groups 4 and 5 can be filled with these, and a phrase such as \texttt{the tan hen sat on the nest} is available for Group~6. The QWERTY first exercise has no word as common as \texttt{the}. Table~\ref{tab:vocab} in Chapter~\ref{ch:constrained} quantifies the difference: 185 admissible words against 21 in the cleaned frequency list, and about a fifth of rank-weighted frequency against about one percent. The generated corpus of Appendix~G shows what the difference looks like on the page. Exercise~1 for Dvorak by position begins with the chains above, then the words \texttt{the that not has}, then \texttt{that these than state}, then \texttt{states estate status season}.

\section{The Cost of Matching by Position}

Position matching breaks down at later stages, because the mapping sends some QWERTY new keys to punctuation.

\begin{center}
\begin{tabular}{@{}clp{6.8cm}@{}}
\toprule
Ex. & Dvorak keys & Effect on the lesson \\
\midrule
1 & A O E U S N T H & rich vocabulary; no semicolon in the lesson \\
2 & I D & completes the home row; sentences become possible \\
6 & comma, C & the comma is a punctuation key, so words cannot exercise it; sentences must \\
13 & semicolon & no letter at all; the lesson can only be sentences with clauses \\
\bottomrule
\end{tabular}
\end{center}

At Exercise~6 the new key W maps to the comma, so a sentence such as \texttt{the students, the tenants, count the stones} is needed to exercise it. At Exercise~13 the Z position holds the semicolon, so the salience regime requires syntactic structure in place of words, since no word contains a semicolon. Matching by position exposes a fact that matching by letter conceals, namely that the difficulty of a key depends on its role in language as well as its location.

\section{Matching by Letter}

The alternative is to keep the QWERTY letter order and let the geometry differ. The vocabulary $\Vocab_i$ is preserved, so the lexicon of each lesson is the same across layouts, but the diagram changes and the learner's movement at each stage differs. The two approaches answer different questions: position matching asks what the same movements permit, and letter matching asks what the same words require.

\section{The Letter Region}

A useful figure for each layout is the outline of the letter region by row.

\begin{center}
\begin{tabular}{@{}lll@{}}
\toprule
Row & QWERTY letters & Dvorak letters \\
\midrule
Top & Q to P (10 keys) & P to L (7 keys) \\
Home & A to L (9 keys) & A to S (10 keys) \\
Bottom & Z to M (7 keys) & Q to Z (9 keys) \\
\bottomrule
\end{tabular}
\end{center}

Dvorak concentrates the letters on the home row and places punctuation where the QWERTY letters were, so the region that a QWERTY typist reaches most often for letters is a region of punctuation in Dvorak.

\section{Transfer Learners and Region Mnemonics}

A person retraining from QWERTY to Dvorak differs from the novice of the paper tutors, because habits already exist and must change. A set of handwritten retraining notes shows how such a learner can structure the task.

\begin{itemize}
\item \textbf{Start from invariants.} The notes begin with the two letters whose positions do not change, A and M, and drill them in combinations: \texttt{a m am ma amma mama}. The learner thereby begins with keys that carry no interference.
\item \textbf{Name the displaced keys.} The notes record the interference directly, for example that the key in the QWERTY F position is U and the key in the J position is H, with the instruction to use the home row.
\item \textbf{Drill by vowel and consonant blocks.} The blocks \texttt{aoeui} and \texttt{dhtns} are typed in isolation before words.
\item \textbf{Use vowel-free words as consonant drills.} \texttt{ssh}, \texttt{shh}, \texttt{nth}, \texttt{brr}, \texttt{hmm}.
\item \textbf{Use rare-key words.} \texttt{kook}, \texttt{yippee}, \texttt{yuppie}, \texttt{exequy}, \texttt{opaque}, \texttt{upkeep} exercise the keys on the outer and lower rows.
\end{itemize}

The notes also record acrostic mnemonics that fix the columns by language.

\begin{center}
\begin{tabular}{@{}lll@{}}
\toprule
Region & Keys (top to bottom) & Mnemonic \\
\midrule
Inner left column & Y, I, X & \emph{Yes, its xmas} \\
Inner right column & F, D, B & \emph{find daily bread} \\
Right edge & L, S, Z & \emph{la zy} (marks corners L and Z; S implicit) \\
Left edge & P, A, Q & none recorded \\
\bottomrule
\end{tabular}
\end{center}

Each phrase fixes a column by the initials of its words read in top-to-bottom order, so the mnemonic is a positional index and not a drill. The right-edge mnemonic marks endpoints only. The gap on the left edge is the natural place to add a further aid. These are self-devised aids of a kind that no paper tutor supplies: the paper tutors teach location by repetition and by the highlighted diagram, and this learner adds a fourth route through language.

\chapter{Two Pedagogical Dialects}

Two layouts generate different early English even when both train toward the same competence.

\section{Parallel Passages}

The QWERTY lines below are transcribed from the Academy pages. The Dvorak lines are constructed for illustration, using only the keys that position matching admits at the same exercise number.

\begin{Verbatim}[label=Exercise 1, labelposition=topline]
QWERTY  ass lad add alf
Dvorak  the tenants hasten out at sunset
\end{Verbatim}

\begin{Verbatim}[label=Exercise 2, labelposition=topline]
QWERTY  a sad lad asks dad;
Dvorak  the students hate it and sit in the sun
\end{Verbatim}

\begin{Verbatim}[label=Exercise 6, labelposition=topline]
QWERTY  Ilse who is a redhead is a selfish girl who likes a wild life;
Dvorak  the students, the tenants, count the stones
\end{Verbatim}

The first Dvorak line is a plausible English clause and the first QWERTY line is a list of near-words. By Exercise~6 the new keys of the Dvorak lesson are a comma and C, so the comma has forced a sentence structure onto the lesson, while the QWERTY sentence is selected only to contain W and I.

\section{Convenient versus Representative Text}

A pedagogically convenient sequence is not necessarily a linguistically representative one. A line built from the home row of QWERTY is easy to generate and lexically odd; a line built from the home row of Dvorak is closer to ordinary prose, because English concentrates its frequent letters there. Yet representativeness is not always a virtue in a lesson. A representative lesson exercises the most frequent transitions and may leave rare keys undertrained until late, while a lesson that is salient for the new key deliberately over-samples it. The two dialects show the tension: the geometry of a keyboard sets which sequences are easy to write, and the author's choice of regime decides whether the lesson follows that ease or corrects for it.

\section{What the Dialects Reveal}

Comparing the dialects makes visible three things that a single layout hides. The first is that the strangeness of early exercises is a property of the pair (layout, lexicon) and not of typing instruction as such. The second is that the difficulty of a key depends on its linguistic role, since a punctuation key cannot be exercised by a word list. The third is that the ladder is independent of the layout, which is why it can serve as the common baseline of the comparison.

\chapter*{Study Questions, Part VII}
\addcontentsline{toc}{chapter}{Study Questions, Part VII}
\begin{enumerate}
\item Write a five-line lesson under the alphabet $\{a,o,e,u,h,t,n,s\}$, then a second under the QWERTY home row, and compare their readability.
\item Use the script of Appendix~D with a word list of your choice and report $|\Vocab_1|$ for both layouts. State the list.
\item Why can the comma not be exercised by a word list? What kind of lesson exercises it?
\item Devise a mnemonic for the left edge (P, A, Q) of the Dvorak letter region, in the manner of the three recorded in Chapter~24.
\item Compare position matching and letter matching. Which question does each answer?
\end{enumerate}
